\begin{description}
\item[(a)] We know $m_t-m_e=-2.5\log(F_t/F_e)$. \hfill(1 point)
Since $F_{\lim}\propto 1/D^2$, $m_t-m_e=5\log(1000/7)$. \hfill(1 point)
Thus $m_t=10.8+6=16.8$. \hfill(1 point)
Using the distance modulus,
\[m=M+5\log\left(\frac{r}{10\,\mathrm{pc}}\right),\qquad\frac r{10}=10^{3.454}.\]
So $r=29\,\mathrm{kpc}$ is the maximum distance to which the student can detect the star using the telescope. \hfill(2 points)
\item[(b)] $m'=m+A=16.8+\alpha r$. \hfill(2 points)
Now $m'=M+5\log(r/(10\,\mathrm{pc}))$, so
\[16.8+0.00005r=-0.5+5\log r-5.\]
% sic: The printed intermediate magnitude equation would give a positive extinction term in the exponent. The official iteration formula instead has the physically appropriate negative term. Both equations are retained.
\hfill(1 point)
Therefore $r=10^{4.46-0.00001r}$. \hfill(2 points)
This equation cannot be solved analytically. Iterate, taking the previous result of $29\,\mathrm{kpc}$ as a good initial approximation. \hfill(2 points)
Then $r=18.73\,\mathrm{kpc}$. \hfill(3 points)
\end{description}
